\scanpage{47}
% archive sheet number: 46

la flèche \unsure{résulterait|résulte} aussi de la suite exacte de \unsure{S-groupes|groupes} sur \(S\)

\diagram{\begin{tikzcd}[ampersand replacement=\&,column sep=small]
e\arrow[r] \& \deleted{G}\unsure{\Gamma|H|\Pi}_1^{\deleted{\unsure{P|\rho}}}\arrow[r] \& \Gamma_1^{\deleted{P}}\arrow[r] \& E_1^{\deleted{P}}\arrow[r] \& e \& \Gamma_1^{\deleted{P}}\simeq E_1^{\deleted{P}}\,\deleted{G}\Gamma_1^{\deleted{P}}
\end{tikzcd}}

qui permet d'écrire

\diagram{\begin{tikzcd}[ampersand replacement=\&,column sep=small,row sep=small]
H^0(S,\deleted{\unsure{\Gamma}}E_1^{\deleted{P}})\arrow[r] \& H^1(S,\deleted{G}\Gamma_1^{\unsure{\circ|P|\rho}})\arrow[r] \& H^1(S,\Gamma_1^{\deleted{P}})\arrow[r]\arrow[d,equal] \& H^1(S,E^P)\\
{} \& {} \& H^1(S,\Gamma) \& {}
\end{tikzcd}}

\deleted{On a :} À priori, l'application envisagée n'est pas injective. \deleted{\unsure{Ainsi|Si}}\inserted{Cependant} l'\unsure{ens|un|ens.} des éléments \deleted{\unsure{qui}} de \(H^1(S,\Gamma_1^\circ)\) qui donnent l'élément neutre de \(H^1(S,\Gamma_1^{\deleted{P}})\) \deleted{\unsure{sont ceux qui|sont ceux que}} est en correspondance biunivoque avec \(H^0(S,E^P)/\deleted{\operatorname{Im}}H^0(S,\Gamma^P)\), est trivial puisque \(\Gamma_1^{\deleted{P}}\) est un produit \(\tfrac12\) direct.

Mais si \(\xi\in H^1(S,\deleted{G}\Gamma_1^\circ)\), on peut utiliser \(\xi\) pour tordre encore la situation, on trouve

\diagram{\begin{tikzcd}[ampersand replacement=\&,column sep=small,row sep=small]
e\arrow[r] \& (\deleted{G}\Gamma_1^{\deleted{P}})^\xi\arrow[r] \& (\Gamma_1^{\deleted{P}})^\xi\arrow[r]\arrow[d,equal] \& \deleted{(}E_1^{\deleted{P}}\deleted{)^\xi}\arrow[r] \& e\\
e\arrow[r] \& \operatorname{Ad}(G)\arrow[r] \& \operatorname{Aut}(G)\arrow[r] \& \operatorname{Ext}(G)\arrow[r] \& e
\end{tikzcd}}

et ici, on est ramené à la question si \deleted{\unsure{les|tout}} automorphisme ext. de \(G\) provient d'un automorphisme \(G\) n'est pas du tout évident quand \(G\) n'est \deleted{\unsure{autre chose que|pas un groupe}} pas une forme extérieure type. [i.e. \deleted{\unsure{on puisse|groupe}} n'ait \unsure{un|pas} de ss-groupe de Borel \(\supset\) Tore maximal].

% slanted
% enclosed in a pen oval
\margin{\uline{exemple} : sur un corps \unsure{à trouver|à Bra|de base} ??}

% slanted
\margin{Faux sur \(\mathbb R\) avec les formes \unsure{diverses|tordues} (\unsure{gpes unitaires|gpes unitaires en|gr. unitaires}): deux formes hermitiennes opposées donnent le \(=\) \unsure{gpe|groupe|gr.} Algébrique sur \(\mathbb R\) mais ne sont pas isom.}
